% Written last. Rules (both venues): 150-250 words; objective -> method -> results -> conclusion;
% no citations; no abbreviations; numbers only via macros; says "single crash" and states the H3 null.
\begin{abstract}
When a software agent that makes payments is stopped, revocation prevents its next actions but not the payments it already created, which stay pending until the merchant captures them. An operator must cancel the suspect ones before capture, with limited capacity and an interface whose acknowledgments can be lost. We define a disposition contract that separates prevented, accounted-for and unresolved payments and lets a payment be reported as cancelled only on evidence from the payment processor. A pre-registered study of \NRunsPrimary{} primary simulator runs, each graded by an independent checker, then measures what that evidence costs. A runbook that trusts acknowledgments reports false closures in every configuration with transport faults, \FCShareAckHigh{} of its cancellation reports under high faults. Verifying each timeout before sending the next cancel costs loss when workers are scarce, whereas sending every cancel first and reconciling afterwards is never significantly worse than trusting acknowledgments and removes every false closure. Re-planning at every event adds nothing measurable. Restarted from a write-ahead journal behind idempotent cancels, evidence-bounded remediation survives a single crash per episode at every instrumented point, with no false closure or duplicate effect. The simulator's transitions match the Stripe test environment, and a case study with a revoked language-model agent shows that the agent's own summary cannot close an incident. Containment is not closure: closure needs evidence, and evidence is cheapest when the remediator does not wait for it.
\end{abstract}

\begin{IEEEkeywords}
reliability engineering, autonomous agents, revocation, incident remediation, runtime assurance
\end{IEEEkeywords}
